% Standalone appendix document - compiles on its own:
%   latexmk -pdf -outdir=build 02-integration.tex
\documentclass[a4paper,11pt]{article}
\usepackage[english]{babel}
\usepackage{../../appendix-style}

\title{5.2 Integration}
\author{SW4PRJ4 -- Group 3}
\date{\today}

\begin{document}
\maketitle

\begin{versionhistory}
  1.0 & 28-10-2026 & Group 3 & Dependency tree inserted \\
  1.1 & 04-11-2026 & Group 3 & Integration strategy and plan added \\
  1.2 & 09-11-2026 & Group 3 & Authentication and authorization plan added \\
  1.3 & 11-11-2026 & Group 3 & Integration overview added \\
\end{versionhistory}

\section{Purpose}
An integration test checks that parts of the system work together through their interfaces. This appendix gives us one strategy for all our integration tests, so we can plan them and build them in a fixed order. We start integrating a class when its unit tests pass.

\section{Dependency Tree}
Figure~\ref{fig:dependency-tree} shows which parts of the system depend on which. It is used to decide the order in which the parts are integrated and tested together. A part can be integrated when the parts below it are in place.

\drawiofig[0.8\linewidth]{dependency-tree}{Dependency tree of the system as it is built today.}{fig:dependency-tree}

At the bottom are the database and the DbContext. The repositories in Infrastructure use the DbContext, the services in Application use the repositories, and the controllers in the API use the services. The two apps call the controllers through the generated api-client.

Keycloak is a separate branch of the tree. The apps log in at Keycloak, and the API's authentication validates the tokens. The controllers do not call Keycloak, but they need the request to be authenticated, so integration tests of the controllers need either Keycloak or test tokens.

\section{Integration Overview}

\subsection{Strategy}
We use collaboration integration. Each column in Table~\ref{tab:integration-overview} is a feature. Each feature gets an integration plan of its own.

\subsection{Rules for All Integration Plans}
\begin{itemize}
  \item A part is integrated only when its unit tests pass.
  \item We are allowed to use a stub instead of Keycloak, since it is outside our system.
  \item The database engine is SQLite in the tests and SQL Server in production. We accept that difference, except in step 2 of the database connection plan (Section~\ref{sec:db-plan}).
  \item A step can have several test cases, for example with a valid, a missing and an expired access token.
  \item In a table representing an integration plan:
    \begin{itemize}
      \item T is the "top module" that the test drives.
      \item X is a module that is included.
      \item S is a stub.
    \end{itemize}
\end{itemize}

\subsection{Database in the Tests}
We use SQLite in the tests as much as possible, because:
\begin{itemize}
  \item The tests are fast.
  \item They need no database server, so they run in CI without extra setup.
  \item Each test can start from an empty database, so the tests do not affect each other.
\end{itemize}

We use a real SQL Server only when SQLite cannot show what we need, for example that the migrations run and that the maximum lengths are enforced (Section~\ref{sec:db-plan}). A real SQL Server can also run in CI as a Docker container, but it is slower and needs more setup.

Table~\ref{tab:sql-server-tests} shows the steps that use a real SQL Server. The tests start the SQL Server container themselves, so Docker must be running on the machine.

\begin{table}[H]
\centering
\small
\rowcolors{2}{gray!15}{white}
\begin{tabular}{@{}l l p{6.5cm}@{}}
\textbf{Integration plan} & \textbf{Step} & \textbf{What it checks} \\
\hline
Database connection & Step 2 & The migrations run, and the maximum lengths are enforced \\
\end{tabular}
\caption{Steps that use a real SQL Server.}
\label{tab:sql-server-tests}
\end{table}

\subsection{Overview of the Integration Plans}
Table~\ref{tab:integration-overview} shows the features in the order we integrate them. Each feature has its own plan and its own subsection in this document. A column shows the parts that are included when on the last step of the feature's integration plan.

The question mark for the system clock means that the choice between a real clock and a stub is made in the subscription plan.
\begin{table}[H]
\centering
\footnotesize
\newcolumntype{C}[1]{>{\centering\arraybackslash}p{#1}}
\rowcolors{2}{gray!15}{white}
\setlength{\tabcolsep}{4pt}
\begin{tabular}{@{}l C{1.7cm} C{2.4cm} C{1.5cm} C{1.8cm} C{1.9cm}@{}}
\textbf{Component} & \textbf{\mbox{Database}\newline \mbox{connection}} & \textbf{\mbox{Authentication}\newline \mbox{and authorization}} & \textbf{\mbox{Get items}} & \textbf{\mbox{Onboarding}} & \textbf{\mbox{Subscription}} \\
\hline
mobile-app                       &   & T &   &   &   \\
admin-panel                      &   & T &   &   &   \\
api-client                       &   & X &   &   &   \\
Controllers                      &   & X & T & T & T \\
AdminOnly policy                 &   & X &   &   &   \\
KeycloakRoleClaimsTransformation &   & X &   &   &   \\
JwtBearer                        &   & X & X & X & X \\
Services                         &   &   & X & X & X \\
Repositories                     & T &   & X & X & X \\
System clock                     &   &   &   &   & ? \\
DbContext                        & X &   & X & X & X \\
Database                         & X &   & X & X & X \\
Keycloak                         &   & X & S & S & S \\
\end{tabular}
\caption{Overview of the integration plans, in the order they are integrated.}
\label{tab:integration-overview}
\end{table}

\section{Integration plan: Database connection}
\label{sec:db-plan}

\begin{table}[H]
\centering
\small
\rowcolors{2}{gray!15}{white}
\begin{tabular}{@{}l c c@{}}
\textbf{Component} & \textbf{Step 1} & \textbf{Step 2} \\
\hline
Repositories & T & T \\
DbContext    & X & X \\
Database     & X (SQLite) & X (SQL Server) \\
\end{tabular}
\caption{Integration plan for the database connection.}
\label{tab:db-integration-plan}
\end{table}

\subsubsection*{Step 1 -- Do the repositories save and read data correctly through the DbContext?}
Since this step focuses on testing that cooperation between repositories and DbContext works properly.
We use SQLite, so the tests run faster.
We are mainly interested in testing to things:
\begin{enumerate}
  \item Mapping between C\# objects and a row in a table.
  \item That the unique KeycloakId works.
\end{enumerate}

\subsubsection*{Step 2 -- Do the migrations and constraints work on SQL Server?}
SQLite does not enforce maximum lengths (which we have some constraints about in e.g. Customer class), and it does not run our migrations.
So for this step, we use a real SQL Sever in Docker, so we can check the constraints that SQLite ignores and run migrations.

\section{Integration plan: API-to-database (DELETE?)}

\subsection{Strategy}
We use collaboration integration. A step in the plan is a feature, and a step integrates the parts that the feature needs. This fits how we work. We build a thin skeleton through all the layers first, with one or a few endpoints in the apps, the API, the database and Keycloak. Then we add features on top of it. Each feature is one more step.

Inside a feature we integrate from the bottom. When a class has another real class below it, we first test the two together, and then the whole feature through the endpoint. For example, the subscription repository is tested against the database before the subscription endpoint is tested. If a test fails, we then know which connection is wrong.

The weak spot of collaboration integration is that some connections in the tree may never be used by a feature. We handle that with a last step that goes through every connection and checks that a test uses it.

\subsection{Rules}
\begin{itemize}
  \item A step drives one top part. The test calls the system the way a user would, through an endpoint or a service.
  \item The parts below the top are real, unless a rule below says otherwise.
  \item Keycloak is always faked with test tokens. It is outside our system, and we do not want the tests to need a running server.
  \item The database engine is SQLite in the tests and SQL Server in production. We accept that difference.
\end{itemize}

\subsection{Plan}
Table~\ref{tab:integration-plan} shows one step for each feature we have today. T is the part the test drives, X is a real part and S is a fake. An empty cell means the part is not used by the feature.

\begin{table}[H]
\centering
\small
\rowcolors{2}{gray!15}{white}
\begin{tabular}{@{}l c c c c@{}}
\textbf{Part} & \textbf{1 Admin access} & \textbf{2 Get items} & \textbf{3 Onboard} & \textbf{4 Subscribe} \\
\hline
Keycloak & S & S & S & S \\
Authentication and roles & X & X & X & X \\
Admin controller & T & & & \\
Items controller & & T & & \\
User controller & & & T & \\
Subscription controller & & & & T \\
Item service & & X & & \\
Customer service & & & X & \\
Subscription service & & & & X \\
Item repository & & X & & \\
Customer repository & & & X & X \\
Subscription repository & & & & X \\
System clock & & & & X \\
DbContext & & X & X & X \\
\end{tabular}
\caption{Integration plan for the features we have today.}
\label{tab:integration-plan}
\end{table}

A last step goes through every connection in the dependency tree and checks that at least one of the steps above uses it. A connection that no step uses becomes a new test.

\subsection{Test Cases}
There is no tool that measures how many connections an integration test covers, so we choose the test cases ourselves. For each step we start from the feature and test what a user can do with it, including the cases where it must fail. We also look at how the unit tests of the lower classes use their interfaces, and we check that the same cases can be sent through the top part.

\subsection{Example: Admin Access}
Step 1 in the plan is the admin endpoint. Only a user with the admin role may call it. A request without a valid token must get the answer unauthorized, and a request with a valid token but without the admin role must get the answer forbidden.

The admin controller is the top part (T). The authentication and the role check are real (X). Keycloak is faked (S). The tests make their own tokens with a test key, and the API is set to trust that key.

The class that reads the roles from the token is already tested alone. The integration test checks the connections around it. They only exist when the API runs, so a test of a single class cannot find an error in them:

\begin{itemize}
  \item The token settings in the API: the key, the issuer, the audience and the expiry.
  \item The order of the middleware, with authentication before authorization.
  \item The registration of the role reader, so the roles from the token reach the policy.
  \item The admin policy, which must give unauthorized without a token and forbidden without the admin role.
\end{itemize}

The test cases and the result are in appendix 5.2.1.

If one of these is wrong, the tests of the single classes still pass, but the endpoint gives the wrong answer. If we forget to register the role reader, an admin gets the answer forbidden.

We build the tests in two stages. Each stage adds one more real part, so an error always points to the part that was added last.

\begin{enumerate}
  \item \emph{The token check.} We send requests where the token is missing, is not a JWT, is signed with another key, has expired, or is made for another audience or issuer. All of them must give unauthorized. The role check is never reached here, so a failure is an error in the token settings or in the order of the middleware.
  \item \emph{The role check.} The token is valid, and only the roles change. A token without the admin role must give forbidden, and a token with the admin role must give ok, also when the user has more roles. The token check is already tested, so a failure is an error in the registration of the role reader or in the policy.
\end{enumerate}

If we only had one test with a valid admin token, a failure could come from any of the four connections, and we would have to search. With the tests in this order, the first failing stage tells us where to look.

\section{Integration plan: Authentication and Authorization}
\label{sec:auth-plan}

\subsection{Scope}
This plan covers authentication and authorization in the API and the clients. The login flow against a real Keycloak is only tested manually in the final step. Onboarding is not part of this plan.

\subsection{Strategy}
The isolated integration plan for testing authentication and authorization follows the strategy ``sandwich''.
\begin{itemize}
  \item Steps 1--3: Bottom-up, from JwtBearer to the API.
  \item Steps 4--5: Top-down, from the apps to the API.
  \item Step 6: Top and bottom meet in the middle.
\end{itemize}

However, the decision to separately test the flow of these components is part of the ``collaboration'' strategy, which applies to the whole project overall.
\subsection{Rules}
These rules apply to the plan in Section~\ref{sec:auth-plan}.
\begin{itemize}
  \item Steps 1--5 are automated and run without a real Keycloak. Keycloak is faked with test tokens, because we do not want the automated tests to need a running server.
  \item A part below the top may be a stub when the real part cannot run in the test, for example JwtBearer in step 1 and the API in steps 4--5.
  \item Step 6 uses a real Keycloak, because that is the only way to check real access tokens.
\end{itemize}

\subsection{Plan}
Table~\ref{tab:auth-integration-plan} shows the steps. T is the top module, X is a real part and S is a stub. An empty cell means the part is not part of the step.

\begin{table}[H]
\centering
\small
\rowcolors{2}{gray!15}{white}
\begin{tabular}{@{}l c c c c c c@{}}
\textbf{Component} & \textbf{Step 1} & \textbf{Step 2} & \textbf{Step 3} & \textbf{Step 4} & \textbf{Step 5} & \textbf{Step 6} \\
\hline
Apps: mobile-app and admin-panel &   &   &   &   & T & T \\
api-client                       &   &   &   & T & X & X \\
Controllers                      &   &   & T & S & S & X \\
AdminController                  &   & T &   & S & S & X \\
AdminOnly                        &   & X &   &   &   & X \\
KeycloakRoleClaimsTransformation & T & X &   &   &   & X \\
JwtBearer                        & S & X & X &   &   & X \\
Keycloak                         &   & S & S &   & S & X \\
\end{tabular}
\caption{Integration plan for authentication and authorization.}
\label{tab:auth-integration-plan}
\end{table}

\subsection{Steps}

\subsubsection*{Step 1 -- Can KeycloakRoleClaimsTransformation correctly translate the realm roles in ClaimsPrincipal?}
Since Transformation is T (top module), JwtBearer must be S (stub), because there is no http-request to the API, which is what would have triggered an automatic call of JwtBearer and its conversion of an access token to a ClaimsPrincipal. Instead of JwtBearer, the tests build the ClaimsPrincipal themselves.

\subsubsection*{Step 2 -- Is the AdminOnly policy correctly enforced based on the access token's validity and roles?}
This time, JwtBearer is X, because the tests send http-requests to AdminController. Keycloak is S (stub), because the tests run without a real Keycloak. The stub for Keycloak consists of two parts:
\begin{itemize}
  \item The RSA key pair: A private key used to sign the access token and a public key that is later used to verify the signature on the access token.
  \item A method that provides signed access tokens.
\end{itemize}

\subsubsection*{Step 3 -- Is access to a non-admin controller correctly enforced based on the access token's validity?}
Non-admin controllers with [authorize] should only allow acces for requests with a valid access token.
So the tests in this step should focus on the \textit{mechanism} behind the enforcement of this rule.
Note 1: The [authorize] functionality should still be tested for each individual controller in the different feature integration plans.
Note 2: KeycloakRoleClaimsTransformation \textit{does} run automatically, since it is registered as IClaimsTransformation and runs (only) when JwtBearer has approved the access token.
However, the result of the transformation does not make a difference in the tests.
What matters for the request to reach the controller is if JwtBearer sets IsAuthenticated to true in the ClaimsPrincipal that it builds from the access token.

\subsubsection*{Step 4 -- When we call methods on api-client, does it attach the access token and handle a 401 response from the API correctly?}
The API is S (stub), and the test driver writes the http response itself.

\subsubsection*{Step 5 -- Do the apps sign in, store the tokens and react correctly to the API's responses?}
Keycloak and the API are S. The apps cannot run against Keycloak in automated tests.

\subsubsection*{Step 6 -- Manual test with all components}
Keycloak is no longer stubbed, nor is any other component. This is the final step where the sandwich meets in the middle. Note that this is the only step that checks real access tokens from Keycloak.

\section{Integration plan: Get Items}
\label{sec:items-plan}

\begin{table}[H]
\centering
\small
\rowcolors{2}{gray!15}{white}
\begin{tabular}{@{}l c c@{}}
\textbf{Component} & \textbf{Step 1} & \textbf{Step 2} \\
\hline
ItemsController & T & T \\
ItemService     & X & X \\
ItemRepository  & X & X \\
JwtBearer       &   & X \\
DbContext       & X & X \\
Database        & X (SQLite) & X (SQLite) \\
Keycloak        &   & S \\
\end{tabular}
\caption{Integration plan for get items.}
\label{tab:items-integration-plan}
\end{table}

ItemService only passes the call on to ItemRepository, so it is integrated together with the controller. ItemRepository against the database is covered by the database connection plan (Section~\ref{sec:db-plan}).

\subsubsection*{Step 1 -- Does ItemsController return the items in the database?}
The tests create ItemsController with the real ItemService and call GetItems \textit{directly}, without an HTTP request.
JwtBearer is therefore not part of the step, and the step only tests the path from the controller to the database.
The tests check that the result has the id and name of each item, and that an empty table gives an empty list.

\subsubsection*{Step 2 -- Is GET /items only available to a logged-in user?}
Now the tests send HTTP requests, so JwtBearer is X and Keycloak is S. A request with a valid test token gets 200 with the items as JSON, and a request without a valid token gets 401. This checks that ItemsController has the [Authorize] attribute, and that the route and the JSON answer work.

\section{What the Plan Does Not Cover}
The tests with SQLite do not test the migrations, because the database is created from the model. They also do not test the apps calling the API. We leave both out of this plan.

\end{document}
